% Generated draft: requires mathematical and source review.
\documentclass[11pt]{article}
\usepackage[margin=25mm]{geometry}
\usepackage{fontspec}
\setmainfont{Latin Modern Roman}
\usepackage{amsmath,amssymb,mathtools,unicode-math}
\setmathfont{Latin Modern Math}
\usepackage{enumitem,needspace,xurl,hyperref}
\setlength{\parindent}{0pt}
\setlength{\parskip}{5pt}
\setlength{\emergencystretch}{3em}
\newcommand{\N}{\mathbb{N}}
\newcommand{\Z}{\mathbb{Z}}
\newcommand{\Q}{\mathbb{Q}}
\newcommand{\R}{\mathbb{R}}
\newcommand{\C}{\mathbb{C}}
\newcommand{\F}{\mathbb{F}}
\newcommand{\A}{\mathbb{A}}
\newcommand{\PP}{\mathbb{P}}
\newcommand{\E}{\mathbb{E}}
\newcommand{\eps}{\varepsilon}
\newcommand{\dd}{\,\mathrm{d}}
\begin{document}
\section*{UnsolvedMath ID: 701310}
\textbf{Source problem code:} catalogue entry
\subsection{Definitions and mathematical statement}
Let \(G\) be a group with identity \(1_G\). For \(x,y\in G\), define the left-normed iterated commutators by
\[
[x,{}_1 y]=x^{-1}y^{-1}xy,\qquad [x,{}_{n+1}y]=[[x,{}_n y],y]\quad(n\geq 1).
\]
An element \(a\in G\) is a \emph{right Engel element} if
\[
\forall g\in G\;\exists n\geq 1:\quad [a,{}_n g]=1_G.
\]
It is a \emph{bounded right Engel element} if
\[
\exists N\geq 1\;\forall g\in G:\quad [a,{}_N g]=1_G.
\]
It is a \emph{left Engel element} if
\[
\forall g\in G\;\exists n\geq 1:\quad [g,{}_n a]=1_G.
\]
Here all iteration indices are positive integers. In the unbounded definitions, the index may depend on \(g\); in the bounded definition, \(N\) is independent of \(g\).

\begin{enumerate}
\item[(a)] Does there exist a group \(G\) and a right Engel element \(a\in G\) such that
\[
\exists h\in G\;\forall n\geq 1:\quad [h,{}_n a]\neq 1_G?
\]
\item[(b)] Does there exist a group \(G\) and a bounded right Engel element \(a\in G\) such that
\[
\exists h\in G\;\forall n\geq 1:\quad [h,{}_n a]\neq 1_G?
\]
\end{enumerate}
No finiteness or other restriction on \(G\) is imposed in the supplied statement. The Engel terminology is made explicit using the commutator convention above. The supplied registry labels the problem open; that label is not an independently verified mathematical or literature assessment.
\subsection{Short English statement}
A. Abdollahi asks whether some group has a right Engel element that is not left Engel, and whether such an example can have a uniform bound on its right Engel iterations. Failure to be left Engel requires one fixed group element whose iterated commutators with the proposed element never become trivial.
\subsection{Sources}
\begin{itemize}
\item Local UnsolvedMath registry, permanent ID \(701310\), exact source code \texttt{catalogue entry}, title \emph{Kourovka Notebook Problem 17.8}. Supplied status: \texttt{open}. Record updated \texttt{2026-10-16T14:00:00Z}; upstream revision \texttt{d371ee917ac9647fbb06cd6fc529deec6f3dda82}.
\item \emph{The Kourovka Notebook (Unsolved Problems in Group Theory)}, original problem 17.8, proposed by A. Abdollahi, 17th issue (2010). The supplied background also identifies the 21st issue edition (2026), \texttt{arXiv:1401.0300v48}. The supplied collection description names E. I. Khukhro and V. D. Mazurov as editors. These bibliographic details and the open label are supplied source metadata, not independently verified here.
\end{itemize}
\end{document}
